\section{Sujet n\textsuperscript{o}\,16}
\noindent\begin{minipage}{0.5\linewidth}\footnotesize Office du Baccalauréat du Cameroun\end{minipage}%
\hfill\begin{minipage}{0.45\linewidth}\footnotesize\raggedleft Examen : \textbf{Baccalauréat} · Série : \textbf{C/E}\\
Épreuve : \textbf{Mathématiques} · Durée : \textbf{4\,h} · Coef. : \textbf{7}\end{minipage}
\par\smallskip{\color{onbo}\rule{\linewidth}{1pt}}

\subsection*{Partie A — Évaluation des ressources \hfill (15 points)}

\noindent\textbf{Exercice 1.} \hfill\textit{(3 points)}\\
La scierie « Bois du Ntem » de Kribi emploie $7$ scieurs et $5$ menuisiers. Pour la
campagne de production, le chef d'atelier doit constituer des équipes de travail.
\begin{enumerate}[label=\arabic*)]
\item De combien de façons peut-il former une équipe de $5$ ouvriers comprenant
\emph{exactement} $3$ scieurs et $2$ menuisiers ? \hfill\textit{(1 pt)}
\item De combien de façons peut-il former une équipe de $5$ ouvriers comprenant
\emph{au moins} $4$ scieurs ? \hfill\textit{(1 pt)}
\item Dans une équipe de $5$ ouvriers, il faut désigner un chef d'équipe, un
adjoint et un magasinier (trois rôles distincts). De combien de façons peut-il
attribuer ces trois rôles à partir des $12$ ouvriers ? \hfill\textit{(1 pt)}
\end{enumerate}

\smallskip
\noindent\textbf{Exercice 2.} \hfill\textit{(3 points)}\\
La scierie dispose de deux grumes (troncs) de longueurs respectives \SI{420}{\centi\metre}
et \SI{264}{\centi\metre}. On souhaite les débiter en planches \emph{toutes de même
longueur}, la plus grande possible, sans aucune perte.
\begin{enumerate}[label=\arabic*)]
\item Déterminer, à l'aide de l'algorithme d'Euclide, le PGCD de $420$ et $264$. \hfill\textit{(1 pt)}
\item En déduire la longueur de chaque planche et le nombre total de planches obtenues. \hfill\textit{(1 pt)}
\item Calculer le PPCM de $420$ et $264$. Deux scies circulaires démarrent
ensemble ; l'une effectue un cycle toutes les $\SI{420}{\second}$, l'autre toutes les
$\SI{264}{\second}$. Au bout de combien de temps redémarreront-elles
simultanément ? \hfill\textit{(1 pt)}
\end{enumerate}

\smallskip
\noindent\textbf{Exercice 3.} \hfill\textit{(4 points)}\\
Le rendement (en unités convenables) d'une découpe de planches est modélisé par la
fonction $f$ définie sur $]0\,;+\infty[$ par $f(x)=x-x\ln x$, de courbe $(\mathcal{C})$.
\begin{enumerate}[label=\arabic*)]
\item Calculer $\displaystyle\lim_{x\to 0^+}f(x)$ et $\displaystyle\lim_{x\to+\infty}f(x)$. \hfill\textit{(1 pt)}
\item Montrer que $f'(x)=-\ln x$ et dresser le tableau de variations de $f$. \hfill\textit{(1,5 pt)}
\item À l'aide d'une intégration par parties, calculer
$\displaystyle I=\int_1^{\mathrm{e}}f(x)\,\dd x$. \hfill\textit{(1,5 pt)}
\end{enumerate}

\smallskip
\noindent\textbf{Exercice 4.} \hfill\textit{(5 points)}\\
L'espace est rapporté à un repère orthonormé direct $(O\,;\,\vec\imath,\vec\jmath,\vec k)$.
On considère les points $A(2\,;0\,;0)$, $B(0\,;3\,;0)$, $C(0\,;0\,;4)$ et $D(2\,;3\,;4)$,
sommets d'une structure en bois.
\begin{enumerate}[label=\arabic*)]
\item Calculer les coordonnées du produit vectoriel $\vec{AB}\wedge\vec{AC}$. \hfill\textit{(1,25 pt)}
\item En déduire une équation cartésienne du plan $(ABC)$. \hfill\textit{(1,25 pt)}
\item Calculer la distance du point $D$ au plan $(ABC)$. \hfill\textit{(1,25 pt)}
\item En déduire le volume du tétraèdre $ABCD$. \hfill\textit{(1,25 pt)}
\end{enumerate}

\subsection*{Partie B — Évaluation des compétences \hfill (5 points)}

\noindent\textit{Monsieur Eyenga dirige une scierie à Kribi. Il sollicite ton aide,
en tant qu'élève de Terminale~C, pour trois décisions. Il dispose de planches
rectangulaires de \SI{6}{\metre} de long et veut y découper des panneaux carrés
d'arête $x$ (en mètres) surmontés d'un demi-disque de diamètre $x$, de sorte que la
hauteur totale d'un panneau ne dépasse pas la largeur disponible ; il cherche à
\emph{maximiser l'aire} d'un tel panneau découpé dans un demi-cercle de rayon
\SI{1}{\metre} (l'aire vaut $S(x)=x\sqrt{1-x^2}$ pour $x\in\,]0\,;1[$, où $2x$ est la
base et $\sqrt{1-x^2}$ la hauteur du rectangle inscrit). Sa production mensuelle de
planches, de $1500$ unités le premier mois, augmente de \SI{6}{\percent} chaque mois.
Enfin, au contrôle qualité, \SI{4}{\percent} des planches présentent un défaut ; un
client en choisit $5$ au hasard.}

\begin{enumerate}[label=\textbf{Tâche \arabic*.},leftmargin=2.4em]
\item Déterminer la valeur de $x$ qui rend l'aire $S(x)=x\sqrt{1-x^2}$ maximale, ainsi
que cette aire maximale. \hfill\textit{(1,5 pt)}
\item Déterminer le mois à partir duquel la production mensuelle dépassera $2500$ planches. \hfill\textit{(1,5 pt)}
\item Déterminer la probabilité qu'\emph{exactement une} des $5$ planches choisies
présente un défaut. \hfill\textit{(1,5 pt)}
\end{enumerate}
\par\smallskip{\footnotesize\itshape Présentation générale : 0,5 point.}

% ════════════════════════════════════════════════════
\corrige{du sujet 16}

\noindent\textbf{Partie A — Exercice 1.}
1) $\binom{7}{3}\binom{5}{2}=35\times10=\mathbf{350}$ équipes.
2) « Au moins $4$ scieurs » sur $5$ ouvriers : soit $4$ scieurs et $1$ menuisier,
soit $5$ scieurs. $\binom{7}{4}\binom{5}{1}+\binom{7}{5}\binom{5}{0}=35\times5+21\times1=175+21=\mathbf{196}$.
3) Trois rôles distincts attribués à $3$ ouvriers choisis parmi $12$ : arrangement
$A_{12}^{3}=12\times11\times10=\mathbf{1320}$.

\noindent\textbf{Exercice 2.}
1) Euclide : $420=1\times264+156$ ; $264=1\times156+108$ ; $156=1\times108+48$ ;
$108=2\times48+12$ ; $48=4\times12+0$. Donc $\mathrm{PGCD}(420,264)=\mathbf{12}$.
2) Chaque planche mesure $\SI{12}{\centi\metre}$ ; nombre de planches :
$\frac{420}{12}+\frac{264}{12}=35+22=\mathbf{57}$ planches.
3) $\mathrm{PPCM}=\dfrac{420\times264}{12}=420\times22=\mathbf{9240}$. Les deux scies
redémarrent ensemble au bout de $\SI{9240}{\second}$, soit $2\,\text{h}\,34\,\text{min}$.

\noindent\textbf{Exercice 3.}
1) $\lim_{x\to0^+}x\ln x=0$ donc $\lim_{x\to0^+}f(x)=0$. En $+\infty$ :
$f(x)=x(1-\ln x)\to-\infty$ (car $1-\ln x\to-\infty$).
2) $f'(x)=1-(\ln x+x\cdot\frac1x)=1-(\ln x+1)=-\ln x$. Ainsi $f'(x)>0$ sur
$]0\,;1[$ et $f'(x)<0$ sur $]1\,;+\infty[$ : $f$ croît puis décroît, maximum
$f(1)=1-1\cdot0=1$.
3) IPP avec $u=1-\ln x$ \dots\ On écrit $f(x)=x-x\ln x$.
$\int_1^{\mathrm e}x\dd x=\big[\tfrac{x^2}{2}\big]_1^{\mathrm e}=\frac{\mathrm e^2-1}{2}$.
Pour $\int_1^{\mathrm e}x\ln x\dd x$, IPP ($u=\ln x,\ v'=x$) :
$\big[\tfrac{x^2}{2}\ln x\big]_1^{\mathrm e}-\int_1^{\mathrm e}\tfrac{x}{2}\dd x
=\tfrac{\mathrm e^2}{2}-\big[\tfrac{x^2}{4}\big]_1^{\mathrm e}
=\tfrac{\mathrm e^2}{2}-\tfrac{\mathrm e^2-1}{4}=\tfrac{\mathrm e^2+1}{4}$.
Donc $I=\frac{\mathrm e^2-1}{2}-\frac{\mathrm e^2+1}{4}=\frac{2\mathrm e^2-2-\mathrm e^2-1}{4}
=\dfrac{\mathrm e^2-3}{4}\approx1{,}10$.

\noindent\textbf{Exercice 4.}
1) $\vec{AB}=(-2\,;3\,;0)$, $\vec{AC}=(-2\,;0\,;4)$.
$\vec{AB}\wedge\vec{AC}=(3\cdot4-0\cdot0\,;\,0\cdot(-2)-(-2)\cdot4\,;\,(-2)\cdot0-3\cdot(-2))
=(12\,;8\,;6)$.
2) Le plan $(ABC)$ a pour vecteur normal $\vec n=(12\,;8\,;6)$ (ou $(6\,;4\,;3)$).
Équation : $6x+4y+3z=d$ ; avec $A(2\,;0\,;0)$ : $d=12$. Donc
$\mathbf{6x+4y+3z-12=0}$.
3) $d(D,(ABC))=\dfrac{|6\cdot2+4\cdot3+3\cdot4-12|}{\sqrt{36+16+9}}
=\dfrac{|12+12+12-12|}{\sqrt{61}}=\dfrac{24}{\sqrt{61}}\approx3{,}07$.
4) Aire de $ABC=\frac12\|\vec{AB}\wedge\vec{AC}\|=\frac12\sqrt{144+64+36}=\frac12\sqrt{244}
=\sqrt{61}$. Volume $=\frac13\times\text{aire}\times d=\frac13\times\sqrt{61}\times\frac{24}{\sqrt{61}}
=\frac{24}{3}=\mathbf{8}$ (unités de volume).

\smallskip
\noindent\textbf{Partie B — Situation.}
\emph{Tâche 1.} $S(x)=x\sqrt{1-x^2}$. $S(x)^2=x^2(1-x^2)=x^2-x^4$ ; posons
$g(x)=x^2-x^4$, $g'(x)=2x-4x^3=2x(1-2x^2)$, nul (sur $]0;1[$) en
$x=\frac{1}{\sqrt2}=\frac{\sqrt2}{2}$. C'est un maximum. Alors
$S\!\left(\tfrac{1}{\sqrt2}\right)=\tfrac{1}{\sqrt2}\sqrt{1-\tfrac12}=\tfrac{1}{\sqrt2}\cdot\tfrac{1}{\sqrt2}=\mathbf{\tfrac12}$.
Donc $x=\frac{\sqrt2}{2}\approx0{,}71$ m et aire maximale $\SI{0,5}{\metre\squared}$.

\emph{Tâche 2.} Production du mois $n$ : $u_n=1500\times1{,}06^{\,n-1}$. On résout
$u_n>2500$, soit $1{,}06^{\,n-1}>\frac{5}{3}$, d'où
$n-1>\frac{\ln(5/3)}{\ln1{,}06}\approx\frac{0{,}5108}{0{,}0583}\approx8{,}77$ : à partir du
\textbf{10\textsuperscript{e} mois} ($u_{10}=1500\times1{,}06^9\approx\mathbf{2535}$ planches).

\emph{Tâche 3.} Loi binomiale $\mathcal{B}(5\,;0{,}04)$.
$P(X{=}1)=\binom{5}{1}(0{,}04)^1(0{,}96)^4=5\times0{,}04\times0{,}8493\approx\mathbf{0{,}170}$,
soit environ \SI{17}{\percent}.
